\documentclass[11pt]{article}
% Shared design for the two minimal papers.
% It follows the layout of the guide "The Riemann Hypothesis and its generalizations"
% (KNT with Fable, 9 November 2026): a designed title block, a short story list at the start,
% coloured headings, boxed statements, idea boxes, and annotated figures.
%
% A paper sets \papertitle, \papersubtitle and \paperrunning before \begin{document}
% and calls \maketitleblock after it.

\usepackage[letterpaper,left=1.12in,right=1.12in,top=1.05in,bottom=1.1in]{geometry}
\usepackage{amsmath,mathtools,amssymb,amsthm}
\usepackage{newpxtext}
\usepackage[varbb]{newpxmath}
\usepackage{fontspec}
\setsansfont{texgyreheros}[Extension=.otf, UprightFont=*-regular, BoldFont=*-bold,
  ItalicFont=*-italic, BoldItalicFont=*-bolditalic, Scale=0.88]
\usepackage{microtype}
\usepackage{parskip}
\usepackage{xcolor}
\usepackage{booktabs,array}
\usepackage{enumitem}
\usepackage{multicol}
\usepackage[titles]{tocloft}
\usepackage{tikz}
\usetikzlibrary{decorations.markings,arrows.meta,positioning,calc}
\usepackage[most]{tcolorbox}
\usepackage{fancyhdr}
\usepackage{titlesec}
\usepackage[font=small,labelfont={bf,sf},labelsep=period,skip=6pt]{caption}
\usepackage{url}

\definecolor{accent}{HTML}{1B4F72}
\definecolor{warm}{HTML}{C0562A}
\definecolor{slate}{HTML}{56636E}
\definecolor{ink}{HTML}{1E1E1E}

\usepackage[colorlinks=true,linkcolor=accent,citecolor=accent,urlcolor=accent]{hyperref}

% Lean identifiers: typewriter, breakable at dots and underscores.
\DeclareUrlCommand\lean{\urlstyle{tt}}

% ---------- layout ----------
\linespread{1.06}
\setlength{\parskip}{0.55em plus 0.1em minus 0.05em}
\setlist{itemsep=0.25em, topsep=0.2em, parsep=0.1em}
\titleformat{\section}{\sffamily\bfseries\Large\color{accent}}{\thesection}{0.7em}{}
\titleformat{\subsection}{\sffamily\bfseries\normalsize\color{ink}}{\thesubsection}{0.6em}{}
\titleformat{\paragraph}[runin]{\sffamily\bfseries\small\color{accent}}{}{0pt}{}[\enspace]
\titlespacing*{\section}{0pt}{22pt}{5pt}
\titlespacing*{\subsection}{0pt}{13pt}{1pt}
\titlespacing*{\paragraph}{0pt}{9pt}{0pt}
\pagestyle{fancy}
\fancyhf{}
\renewcommand{\headrulewidth}{0pt}
\fancyhead[L]{\sffamily\footnotesize\color{slate}\paperrunning}
\fancyhead[R]{\sffamily\footnotesize\color{slate}\thepage}
\setlength{\headheight}{14pt}
\setcounter{tocdepth}{1}
\setlength{\cftbeforesecskip}{2pt}
\renewcommand{\cftsecfont}{\small}
\renewcommand{\cftsecpagefont}{\small\color{slate}}
\renewcommand{\cftsecleader}{\cftdotfill{\cftdotsep}}
\setlength{\cftsecnumwidth}{1.7em}
\setlength{\emergencystretch}{2em}
\numberwithin{equation}{section}

% ---------- title block ----------
\newcommand{\papertitle}{}
\newcommand{\papersubtitle}{}
\newcommand{\paperrunning}{}
\newcommand{\paperauthors}{KNT\enspace\textperiodcentered\enspace Astra\enspace\textperiodcentered\enspace Fable~5.2}
\newcommand{\paperdate}{Draft of 11 November 2026}
\newcommand{\maketitleblock}{%
  \thispagestyle{plain}%
  {\sffamily\bfseries\color{accent}\fontsize{25}{29}\selectfont\raggedright\papertitle\par}%
  \vspace{5pt}%
  {\sffamily\Large\color{slate}\raggedright\papersubtitle\par}%
  \vspace{8pt}%
  {\color{accent}\hrule height 1.2pt}%
  \vspace{7pt}%
  {\sffamily\small\color{ink}\paperauthors\hfill\color{slate}\paperdate\par}%
  \vspace{10pt}}

% Small headings inside the front matter.
\newcommand{\fronthead}[1]{\par\vspace{6pt}{\sffamily\bfseries\normalsize\color{ink}#1}\par\vspace{-2pt}}
\makeatletter
\newcommand{\contentsblock}{%
  \par\noindent\begin{minipage}{\linewidth}%
  \fronthead{Contents}%
  \setlength{\parskip}{0pt}\setlength{\columnsep}{2.2em}\raggedright
  \begin{multicols}{2}
  \@starttoc{toc}
  \end{multicols}
  \end{minipage}\par}
\makeatother

% ---------- statements ----------
% Theorems, lemmas, propositions and corollaries share one counter.
% Each statement is in a tinted box with an accent rule at the left.
\newtheoremstyle{boxed}{0pt}{0pt}{\normalfont}{0pt}{\sffamily\bfseries\small\color{accent}}{.}{0.6em}{}
\theoremstyle{boxed}
\newtheorem{theorem}{Theorem}[section]
\newtheorem{lemma}[theorem]{Lemma}
\newtheorem{proposition}[theorem]{Proposition}
\newtheorem{corollary}[theorem]{Corollary}

\tcbset{guidebox/.style={enhanced, parbox=false, boxrule=0pt, frame hidden,
  arc=1.5pt, left=11pt, right=10pt, top=7pt, bottom=7pt, before skip=11pt, after skip=9pt}}
% Not breakable: a label at the start of a box that breaks at the top of a page can be lost.
\tcbset{statementbox/.style={guidebox, colback=accent!6,
  borderline west={3pt}{0pt}{accent}}}
\tcolorboxenvironment{theorem}{statementbox}
\tcolorboxenvironment{lemma}{statementbox}
\tcolorboxenvironment{proposition}{statementbox}
\tcolorboxenvironment{corollary}{statementbox}

% The main result of a paper: a framed box with a title bar.
% \begin{maintheorem}{name} uses the counter of the theorems, and \label works after it.
\newtcolorbox{mainresult}[1]{enhanced, parbox=false, colback=white, colframe=accent, boxrule=0.9pt,
  arc=1.5pt, left=11pt, right=11pt, top=8pt, bottom=8pt, title={#1},
  fonttitle=\sffamily\bfseries\small, colbacktitle=accent, coltitle=white,
  before skip=12pt, after skip=12pt}
\newenvironment{maintheorem}[1]%
  {\refstepcounter{theorem}\begin{mainresult}{Theorem~\thetheorem\ (#1)}}%
  {\end{mainresult}}

% ---------- proofs ----------
% A proof has a thin rule at the left, so that the eye finds the statements.
\renewcommand{\proofname}{\upshape\sffamily\bfseries\small\color{slate}Proof}
\renewcommand{\qedsymbol}{\textcolor{slate}{\rule{0.55em}{0.55em}}}
\tcolorboxenvironment{proof}{enhanced, breakable, parbox=false, boxrule=0pt, frame hidden,
  colback=white, borderline west={1.2pt}{0pt}{slate!35}, left=11pt, right=0pt, top=1pt, bottom=2pt,
  before skip=6pt, after skip=11pt}
% A step or a case inside a proof.
\newcommand{\step}[1]{{\sffamily\bfseries\small\color{slate}#1}}

% ---------- idea, notation and fine-print boxes ----------
\newtcolorbox{idea}[1]{guidebox, breakable, colback=warm!6, borderline west={3pt}{0pt}{warm!80},
  before upper={{\sffamily\bfseries\footnotesize\color{warm!75!black}IDEA\enspace\textbar\enspace}{\sffamily\bfseries\small #1}\par}}
\newtcolorbox{notation}[1]{guidebox, breakable, colback=black!4, borderline west={3pt}{0pt}{slate!60},
  fontupper=\small,
  before upper={{\sffamily\bfseries\footnotesize\color{slate}NOTATION\enspace\textbar\enspace}{\sffamily\bfseries\small #1}\par}}
\newtcolorbox{fineprint}{guidebox, breakable, colback=white, borderline west={1.5pt}{0pt}{black!25},
  fontupper=\footnotesize\color{black!78}, left=9pt, top=1pt, bottom=1pt,
  before upper={{\sffamily\bfseries Fine print.}\enspace}}

% ---------- figures ----------
\tikzset{
  vtx/.style={circle, fill=ink, inner sep=1.35pt},
  port/.style={circle, draw=accent, fill=white, line width=0.8pt, inner sep=1.5pt},
  edge/.style={line width=0.55pt, draw=ink},
  hot/.style={line width=1.1pt, draw=warm},
  cool/.style={line width=1.1pt, draw=accent},
  region/.style={draw=accent, fill=accent!7, line width=0.6pt, rounded corners=6pt},
  note/.style={font=\sffamily\footnotesize, text=slate},
}

% ---------- macros of the mathematics ----------
\newcommand{\Z}{\mathcal{Z}}
\newcommand{\PP}{\mathbb{P}}
\newcommand{\EE}{\mathbb{E}}
\newcommand{\comp}{\operatorname{comp}}
\newcommand{\tr}{\operatorname{tr}}
\newcommand{\lgk}[1]{\log^{(#1)}}
\newcommand{\lgs}{\log^{*}}
\newcommand{\sm}{\setminus}
\newcommand{\half}{\tfrac12}
\newcommand{\cO}{\mathcal{O}}


\renewcommand{\papertitle}{The minimum number of edges\\ in a pancyclic graph}
\renewcommand{\papersubtitle}{A short proof for Erd\H{o}s Problem \#1016}
\renewcommand{\paperrunning}{The minimum number of edges in a pancyclic graph}
\hypersetup{pdftitle={The minimum number of edges in a pancyclic graph: a short proof},pdfauthor={KNT, Astra, Fable 5.2}}

\begin{document}
\maketitleblock

A graph on $n$ vertices is pancyclic if it has a cycle of every length from $3$ to $n$. Let
$n+h(n)$ be the least number of edges of such a graph. We prove that
\[
  \log_2n+\log^{*}n-O(1)\;\le\;h(n)\;\le\;\lfloor\log_2n\rfloor+\log^{*}n+1 .
\]
This answers Problem~\#1016 of Erd\H{o}s. The lower bound follows from one estimate for random
even subgraphs, which is proved in~\cite{Max}. The upper bound is a construction. Both bounds are
formalised in Lean~4. Astra and Fable~5.2 are language models. Section~\ref{sec:status} states
the part of each author.

\fronthead{The proof in five sentences}
\begin{enumerate}[leftmargin=1.6em]
\item A pancyclic graph with $n$ vertices and $n+h$ edges has rank $h+1$, and its cycles have
$n-2$ different lengths. A graph of rank $r$ has fewer than $2^{r}$ cycles. This gives only
$h\ge\log_2n-1$.
\item Take a small subgraph $W$ that already contains cycles of all short lengths. A cycle through
an edge of $W$ is determined by its part outside $W$, up to at most $|E(W)|$ lengths, and this
part is a forest.
\item By the \emph{forest estimate}, at most $\half+o(1)$ of all outside parts are forests. So
the largest possible density of cycle lengths decreases by a factor of almost two when the rank
grows from $R$ to $2^{256R}$.
\item The halving repeats at each scale of a tower of exponentials. This gives the term
$\log^{*}n$ of the lower bound.
\item For the upper bound, $k$ chords with path lengths $1,2,4,\dots$ give $2^{k}$ consecutive
cycle lengths, and $\log^{*}k$ further chords supply the short lengths.
\end{enumerate}

\section{The result}

A simple graph on $n$ vertices is \emph{pancyclic} if it contains a cycle of every length
$3,4,\dots,n$. Write $n+h(n)$ for the least number of edges of such a graph. Throughout, $\log$ is
the binary logarithm, and $\lgk j$ is its $j$-fold iterate. The number $\lgs x$ is the least
nonnegative integer $j$ with $\lgk jx\le1$.

\begin{maintheorem}{the minimum excess of a pancyclic graph}\label{thm:main}
There is a constant $A$ such that, for every $n\ge3$,
\[
  \log n+\lgs n-A\;\le\;h(n)\;\le\;\lfloor\log n\rfloor+\lgs n+1 .
\]
\end{maintheorem}

\begin{center}\small
\begin{tabular}{@{}p{5.4cm}p{3.9cm}p{4.3cm}@{}}
\toprule
 & \textsf{\textbf{lower bound for $h(n)$}} & \textsf{\textbf{upper bound for $h(n)$}}\\
\midrule
Bondy~\cite{Bondy}, stated without proof & $\log(n-1)-1$ & $\log n+\lgs n+O(1)$\\
Theorem~\ref{thm:main} & $\log n+\lgs n-A$ & $\lfloor\log n\rfloor+\lgs n+1$\\
\bottomrule
\end{tabular}
\end{center}

A count of cycles gives the lower bound of Bondy: a connected graph with $n+h$ edges has at most
$2^{h+1}-1$ cycles, and a pancyclic graph needs $n-2$ of them. Griffin~\cite{Griffin} gave a
written proof and computed small values. Problem~\#1016 of Erd\H{o}s~\cite{Bloom} asks if the
lower bound can increase to $\log n+\lgs n-O(1)$, and if at least $h(n)-\log n\to\infty$.
Theorem~\ref{thm:main} answers both questions.

The theorem and its proof are from~\cite{KNT}. This paper gives the proof in a short form. It
uses one result of~\cite{Max}, the forest estimate, which we state in Section~\ref{sec:input}.

\section{The forest estimate}\label{sec:input}

\begin{notation}{even subgraphs and outside parts}
\begin{itemize}[leftmargin=1.2em]
\item For a graph $G$, an \emph{even subgraph} is a set of edges that meets every vertex an even
number of times. With symmetric difference as addition, the even subgraphs form a vector space
$\Z(G)$ over $\mathbb{F}_2$.
\item Its dimension is the \emph{rank} $\beta(G)=|E(G)|-|V(G)|+\comp(G)$, where $\comp(G)$ is the
number of components. A cycle is a nonzero element of $\Z(G)$, so a graph of rank $r$ has at most
$2^{r}-1$ cycles.
\item Let $W$ be a subgraph of $G$, with $r=\beta(G)$, $t=\beta(W)$ and $m=|E(W)|$. For an even
subgraph $\xi$ we call $\xi\sm E(W)$ its \emph{outside part}. The map $\xi\mapsto\xi\sm E(W)$ is
linear on $\Z(G)$ with kernel $\Z(W)$. So there are $2^{r-t}$ outside parts.
\item $p(G,W)$ is the proportion of the outside parts that are forests.
\end{itemize}
\end{notation}

\begin{theorem}[Forest estimate, {\cite[Theorem~3.1]{Max}}]\label{thm:forest}
There is $r_0$ such that the following holds for all $r\ge r_0$. Let $G$ be a connected simple graph
of rank $r$. Let $W\subseteq G$ be a nonempty union of cycles with at most $r/\log^{2}r$ edges and
at most $(\log r)/100$ components. Then
\[
  p(G,W)\;\le\;\frac12+\frac{20}{\lgk{3}r}.
\]
\end{theorem}

\begin{lemma}[Lengths]\label{lem:count}
The cycles of $G$ have at most $2^{r-t}\bigl(p\,m+1-p\bigr)$ different lengths, where $p=p(G,W)$.
\end{lemma}

\begin{proof}
Let $x$ be an outside part, and let $C$ be a cycle with the outside part $x$. If $x$ is not a
forest, then $x\subseteq C$ contains a cycle, so $C=x$. Such an $x$ accounts for at most one
length. If $x$ is a forest, then $C$ has an edge in $W$, and the length of $C$ is
$|x|+|C\cap E(W)|$ with $1\le|C\cap E(W)|\le m$. Such an $x$ accounts for at most $m$ lengths.
There are $p\,2^{r-t}$ outside parts of the second kind and $(1-p)2^{r-t}$ of the first.
\end{proof}

\section{The lower bound}\label{sec:lower}

\begin{idea}{outline of the proof}
A number $\Phi(R)\le1$ measures how well a graph of rank at least $R$ can fill an interval of cycle
lengths. In a graph with many cycle lengths we find a small subgraph $W$ that already contains the
short cycles. By the forest estimate for $W$, the number $\Phi$ decreases by a factor of almost
two when the rank grows from $R$ to $2^{256R}$ (Figure~\ref{fig:tower}).
\end{idea}

\begin{figure}[ht]
\centering
\begin{tikzpicture}[x=1cm,y=1cm]
  \draw[slate!80, line width=0.5pt, -{Stealth[length=4pt]}] (-0.3,0) -- (11.2,0);
  \node[note, anchor=west] at (11.25,0) {rank (not to scale)};
  \foreach \j/\x/\hh in {0/0.8/2.4,1/3.4/1.2,2/6.0/0.6,3/8.6/0.3}{
    \fill[accent!18] (\x-0.33,0) rectangle (\x+0.33,\hh);
    \draw[cool] (\x-0.33,\hh) -- (\x+0.33,\hh);
    \draw[slate!80, line width=0.5pt] (\x,0) -- (\x,-0.1);
    \node[note, text=accent, above=1pt] at (\x,\hh) {$\Phi(R_{\j})$};
  }
  \node[note, below=3pt] at (0.8,0) {$R_0$};
  \node[note, below=3pt] at (3.4,0) {$R_1=2^{256R_0}$};
  \node[note, below=3pt] at (6.0,0) {$R_2=2^{256R_1}$};
  \node[note, below=3pt] at (8.6,0) {$R_3=2^{256R_2}$};
  \foreach \a/\b/\y in {1.25/2.95/1.9,3.85/5.55/1.0,6.45/8.15/0.55}
    \draw[warm, line width=0.6pt, -{Stealth[length=4pt]}] (\a,\y) -- node[note, text=warm, above=-1pt] {\scriptsize almost $\times\frac12$} (\b,\y-0.25);
\end{tikzpicture}
\caption{The density $\Phi$ of cycle lengths at the scales $R_{j+1}=2^{256R_j}$. Each step almost
halves it (Proposition~\ref{prop:halving}). The number of steps below a rank $r$ is $\lgs r-O(1)$.}
\label{fig:tower}
\end{figure}

For an integer $R\ge1$ let
\[
  \Phi(R)=\sup\,\frac{L-2}{2^{\beta(G)}} ,
\]
where the supremum is over all simple graphs $G$ with $\beta(G)\ge R$ and all integers $L\ge2$ such
that $G$ contains a cycle of each length $3,\dots,L$. A graph of rank $\beta$ has at most
$2^{\beta}-1$ cycles, so $\Phi(R)\le1$. The function $\Phi$ is non-increasing.

\begin{lemma}[Witness]\label{lem:witness}
Let $R\ge1$ and $\varphi=\Phi(R)$, and let $G$ contain a cycle of each length $3,\dots,L$, where
$L\ge2^{2R}$. Then $G$ has a subgraph $W$ that is a union of at most $2R$ cycles, with
$\beta(W)\ge2R$ and $|E(W)|\le\varphi\,2^{2R}+2^{R}+5R$.
\end{lemma}

\begin{proof}
Put $W_0=\emptyset$ and $W_{i+1}=W_i\cup C_i$, where $C_i$ is a cycle of $G$ whose length $\ell_i$ is
the least integer $\ge3$ that is not the length of a cycle of $W_i$. Stop as soon as the rank is at
least $2R$. Consider a step, and let $s=\beta(W_i)<2R$. The graph $W_i$ contains cycles of the
$\ell_i-3$ lengths $3,\dots,\ell_i-1$. Therefore $\ell_i-3\le2^{s}-1$. This gives
$\ell_i\le2^{2R-1}+2\le L$, so $C_i$ exists. If $s\ge R$, then also $\ell_i-3\le\varphi\,2^{s}$, by
the definition of $\Phi$. The cycle $C_i$ is not a subgraph of $W_i$, so the rank increases at every
step. So there are at most $2R$ steps, their ranks $s$ are distinct, and
\[
  |E(W)|\;\le\;\sum_i\ell_i\;\le\;\sum_{s=0}^{R-1}\bigl(2^{s}+2\bigr)+\sum_{s=R}^{2R-1}\bigl(\varphi\,2^{s}+3\bigr)
  \;\le\;\varphi\,2^{2R}+2^{R}+5R .\qedhere
\]
\end{proof}

\begin{proposition}[Near-halving]\label{prop:halving}
For every sufficiently large integer $R$,
\[
  \Phi\bigl(2^{256R}\bigr)\;\le\;\Bigl(\frac12+\frac{20}{\lgk{2}R}\Bigr)\Phi(R)+2^{2-R}.
\]
\end{proposition}

\begin{proof}
Let $G$ be a simple graph of rank $r\ge2^{256R}$ with a cycle of each length $3,\dots,L$. We bound
$(L-2)\,2^{-r}$. Identifying a vertex of one component with a vertex of another changes neither the
rank nor the set of cycles, so we may assume that $G$ is connected. If $L<2^{2R}$, then
$(L-2)\,2^{-r}<2^{2-R}$, and we are done.

Otherwise let $W$ be as in Lemma~\ref{lem:witness}, with $m$ edges and rank $t\ge2R$, and let
$\varphi=\Phi(R)$. Then $m\le2^{2R+1}\le\sqrt r\le r/\log^{2}r$ and
$\comp(W)\le2R\le(\log r)/100$. So Theorem~\ref{thm:forest} gives
\[
  p(G,W)\;\le\;\frac12+\frac{20}{\lgk3r}\;\le\;\frac12+\delta,
  \qquad\text{where }\delta=\frac{20}{\lgk2R}.
\]
The second inequality holds because $\lgk3r\ge\log(8+\log R)\ge\lgk2R$. The graph $G$ has cycles of
at least $L-2$ lengths. So Lemma~\ref{lem:count}, with $p=p(G,W)$, and the bounds on $t$ and $m$
give
\[
  \frac{L-2}{2^{r}}\;\le\;\frac{pm+1}{2^{t}}
  \;\le\;\Bigl(\frac12+\delta\Bigr)\bigl(\varphi+2^{-R}+5R\,2^{-2R}\bigr)+2^{-2R}
  \;\le\;\Bigl(\frac12+\delta\Bigr)\varphi+2^{2-R}.\qedhere
\]
\end{proof}

\begin{corollary}\label{cor:phi}
There is a constant $c$ such that $\Phi(r)\le c\,2^{-\lgs r}$ for all $r\ge1$.
\end{corollary}

\begin{proof}
Fix $R_0\ge512$ so large that Proposition~\ref{prop:halving} holds for all $R\ge R_0$. Let
$R_{j+1}=2^{256R_j}$ and $\delta_j=20/\lgk2R_j$. By Proposition~\ref{prop:halving}, the numbers
$a_j=2^{j}\Phi(R_j)$ satisfy
\[
  a_{j+1}\;\le\;(1+2\delta_j)\,a_j+2^{\,j+3-R_j}.
\]
For $j\ge2$ we have $\lgk2R_j=8+256R_{j-2}$, so $\sum_j\delta_j<\infty$. Also
$\sum_j2^{\,j-R_j}<\infty$. Hence the numbers $a_j$ are bounded, say $\Phi(R_j)\le c_1 2^{-j}$.

The numbers $R_j$ grow like a tower of exponentials. To bound $\lgs R_j$, we compare them with an
exact tower. Let $U_0=R_0^2$ and $U_{j+1}=2^{U_j}$, so that $\lgs U_j=j+\lgs U_0$. Then
$R_j\le\sqrt{U_j}$ for all $j$. Indeed, if this holds for $j$, then
\[
  R_{j+1}=2^{256R_j}\;\le\;2^{256\sqrt{U_j}}\;\le\;2^{U_j/2}=\sqrt{U_{j+1}},
\]
because $\sqrt{U_j}\ge R_0\ge512$. Hence $\lgs R_j\le j+\lgs U_0$.

Finally let $r\ge R_0$, and let $j$ be largest with $R_j\le r$. Then
$\Phi(r)\le\Phi(R_j)\le c_12^{-j}$ and $\lgs r\le\lgs R_{j+1}\le j+1+\lgs U_0$. Together these give
$\Phi(r)\le c\,2^{-\lgs r}$ with $c=c_1\,2^{1+\lgs U_0}$. For $r<R_0$ use $\Phi(r)\le1$, and enlarge
$c$ if necessary.
\end{proof}

\begin{proof}[Proof of the lower bound in Theorem~\ref{thm:main}]
Let $G$ be a pancyclic graph with $n$ vertices and $n+h$ edges. It is connected, so its rank is
$r=h+1$. By the definition of $\Phi$ and by Corollary~\ref{cor:phi},
\[
  n-2\;\le\;\Phi(r)\,2^{r}\;\le\;c\,2^{\,r-\lgs r}.
\]
The trivial bound $n-2\le2^{r}-1$ gives $n\le2^{r+1}\le2^{2^{r}}$, hence $\lgs n\le\lgs r+2$. Now
take logarithms in the display:
\[
  h=r-1\;\ge\;\log(n-2)+\lgs r-\log c-1\;\ge\;\log n+\lgs n-\log c-5.\qedhere
\]
\end{proof}

\section{The construction}\label{sec:upper}

\begin{idea}{outline of the construction}
The construction is a chain of $k$ segments of one long cycle. Segment $i$ has a chord between
its two ends. So a path can pass it in two ways, and the two ways differ in length by $2^{i}$.
The chain therefore has $2^{k}$ paths of consecutive lengths, and it uses $k$ chords. A few more
chords close shorter and shorter initial parts of the chain, and these give the short cycles.
The construction follows~\cite{KNT} (Figure~\ref{fig:necklace}).
\end{idea}

\begin{proposition}\label{prop:upper}
For every $n\ge3$ there is a pancyclic graph with $n$ vertices and at most
$n+\lfloor\log n\rfloor+\lgs n+1$ edges.
\end{proposition}

\begin{proof}
For $3\le n\le6$ take a cycle $u_0u_1\cdots u_{n-1}$, with the chord $u_0u_2$ if $n\ge4$ and also
the chord $u_0u_3$ if $n=6$. Now let $n\ge7$.

\step{The necklace.} Let $k\ge2$ be the largest integer with $2^{k}+k+1\le n$, so that
$n\le2^{k+1}+k+1$. Let $H$ be a cycle of length $n$, and let $Q\subseteq H$ be a path that is
composed of consecutive segments $Q_0,\dots,Q_{k-1}$, where $Q_i$ has $2^{i}+1$ edges and joins
$v_i$ to $v_{i+1}$. Thus $Q$ has $2^{k}+k-1\le n-2$ edges. Add the chords $e_i=v_iv_{i+1}$ for
$0\le i<k$. For $1\le j\le k$ let $N_j$ be the union of $Q_0,\dots,Q_{j-1}$ and
$e_0,\dots,e_{j-1}$ (Figure~\ref{fig:necklace}).

A path from $v_0$ to $v_j$ in $N_j$ uses either $Q_i$ or $e_i$ for each $i<j$, and it is shorter by
$2^{i}$ if it uses $e_i$. So $N_j$ has $2^{j}$ such paths, exactly one of each length from $j$ to
$2^{j}+j-1$. The rest $H-Q$ of the outer cycle is a path from $v_k$ to $v_0$ with $n-2^{k}-k+1$
edges. Together with $N_k$ it contains a cycle of each length in the interval
$[\,n-2^{k}+1,\ n\,]$.

\step{Closing chords.} For $2\le j\le k$ the chord $f_j=v_0v_j$ closes the paths of $N_j$ to cycles
of all lengths in the interval $I_j=[\,j+1,\ 2^{j}+j\,]$. Add the chord $f_j$ for $j=k$, then for
$j=\lceil\log k\rceil$, and so on: as long as the last value of $j$ is larger than $2$, the next
value is $\lceil\log j\rceil$. The last value is $2$. If $j'=\lceil\log j\rceil$, then
$2^{j'}\ge j$, so the intervals $I_{j'}$ and $I_j$ overlap or are adjacent. Since $I_2=[3,6]$, these
chords give a cycle of each length in $[\,3,\ 2^{k}+k\,]$.

\step{All lengths.} The intervals $[\,3,\ 2^{k}+k\,]$ and $[\,n-2^{k}+1,\ n\,]$ cover all lengths
from $3$ to $n$, except for the length $2^{k}+k+1$ in the case $n=2^{k+1}+k+1$. In that case add
one more chord, which joins $v_k$ to the neighbour of $v_0$ on $H-Q$. This chord, the edge from that
neighbour to $v_0$, and the path $Q$ form a cycle of length $2^{k}+k+1$.

\step{The number of chords.} For an integer $j\ge3$ we have $\lgs\lceil\log j\rceil=\lgs j-1$, and
$\lgs2=1$. So exactly $\lgs k$ chords $f_j$ were added. The number of all chords is therefore at
most $k+\lgs k+1\le\lfloor\log n\rfloor+\lgs n+1$.
\end{proof}

\begin{figure}[ht]
\centering
\begin{tikzpicture}[
  iv/.style 2 args={postaction={decorate,decoration={markings,
      mark=between positions #1 and #2 step #1 with {\fill (0,0) circle (1.1pt);}}}}]
  \coordinate (v0) at (0,0);
  \coordinate (v1) at (1.8,0);
  \coordinate (v2) at (4.0,0);
  \coordinate (v3) at (6.8,0);
  \coordinate (v4) at (10.8,0);
  \draw[iv={0.5}{0.51}]    (v0) to[bend left=50] node[pos=0.5,above=-1pt] {\scriptsize$Q_0$} (v1);
  \draw[iv={0.3333}{0.67}] (v1) to[bend left=50] node[pos=0.5,above=-1pt] {\scriptsize$Q_1$} (v2);
  \draw[iv={0.2}{0.81}]    (v2) to[bend left=50] node[pos=0.5,above=-1pt] {\scriptsize$Q_2$} (v3);
  \draw[iv={0.1111}{0.89}] (v3) to[bend left=50] node[pos=0.5,above=-1pt] {\scriptsize$Q_3$} (v4);
  \draw[iv={0.2}{0.81}]    (v0) .. controls (0.3,3.2) and (10.5,3.2) .. node[pos=0.5,above=-1pt] {\scriptsize$H-Q$} (v4);
  \foreach \a/\b/\i in {v0/v1/0,v1/v2/1,v2/v3/2,v3/v4/3}
    \draw[cool] (\a) -- node[pos=0.5,above=-2pt,text=accent] {\scriptsize$e_{\i}$} (\b);
  \draw[hot,dash pattern=on 3pt off 2pt] (v0) .. controls (0.9,-1.0) and (3.3,-1.0) .. node[pos=0.5,below=-1pt,text=warm] {\scriptsize$f_2$} (v2);
  \draw[hot,dash pattern=on 3pt off 2pt] (v0) .. controls (0.5,-2.0) and (10.3,-2.0) .. node[pos=0.5,below=-1pt,text=warm] {\scriptsize$f_4$} (v4);
  \foreach \p in {v0,v1,v2,v3,v4} \fill (\p) circle (1.9pt);
  \node[left=3pt] at (v0) {$v_0$};
  \node[below=3pt] at (v1) {$v_1$};
  \node[anchor=north west,inner sep=3pt] at (v2) {$v_2$};
  \node[below=3pt] at (v3) {$v_3$};
  \node[right=3pt] at (v4) {$v_4$};
\end{tikzpicture}
\caption{The construction for $n=24$, where $k=4$. The segments $Q_0,\dots,Q_3$ of the outer cycle
$H$ have $2,3,5,9$ edges and are doubled by the chords $e_0,\dots,e_3$. The chord $f_4$ gives the
cycle lengths $5,\dots,20$, the rest $H-Q$ of the outer cycle gives $9,\dots,24$, and $f_2$ gives
$3,\dots,6$.}
\label{fig:necklace}
\end{figure}

\section{Authorship and verification}\label{sec:status}

KNT is a person. Astra and Fable~5.2 are language models: GPT-6 Astra and Claude Fable~5.2. KNT
obtained Theorem~\ref{thm:main}, the forest estimate and their proofs in September 2026 with
Astra~\cite{KNT}, and directed the work at each step. Fable~5.2 wrote this short form.

The theorem is formalised in Lean~4 with Mathlib, in the repository~\cite{Repo} at commit
\texttt{a5cb5fa}.

\begin{center}\small
\begin{tabular}{@{}p{6.6cm}p{7.4cm}@{}}
\toprule
\textsf{\textbf{Statement}} & \textsf{\textbf{Lean declaration}}\\
\midrule
Both bounds of Theorem~\ref{thm:main}, with absolute additive constants & \lean{Erdos1016.mainTheorem}\\
The upper bound $\lfloor\log n\rfloor+\lgs n+1$ & \lean{Erdos1016.Problem1016.upper_bound_integer}\\
Theorem~\ref{thm:forest} for $|E(W)|\le\sqrt r$ & \lean{Erdos1016.ShortProof.fewComponentForestEstimate}\\
\bottomrule
\end{tabular}
\end{center}

Section~\ref{sec:lower} uses the forest estimate only for $|E(W)|\le\sqrt r$. The Lean kernel
accepts the proofs with the axioms \lean{propext}, \lean{Classical.choice} and \lean{Quot.sound}
only. The kernel check covers these formal statements. No machine has checked the text of this
paper, and no referee has read it.

\begin{thebibliography}{9}\small

\bibitem{Bloom} T.~F.~Bloom, Erd\H{o}s Problem \#1016, \url{https://www.erdosproblems.com/1016},
accessed November 2026.

\bibitem{Bondy} J.~A.~Bondy, Pancyclic graphs~I, \emph{J.~Combin.\ Theory Ser.~B} \textbf{11}
(1971), 80--84.

\bibitem{Griffin} S.~Griffin, Minimal pancyclicity, arXiv:1312.0274 (2013).

\bibitem{KNT} KNT, The minimum number of edges in a pancyclic graph, manuscript, September 2026,
\url{https://jehlp.net/erdos1016/}.

\bibitem{Max} KNT, Astra and Fable~5.2, The maximum number of cycles in a graph of large rank,
draft, November 2026.

\bibitem{Repo} KNT, Astra and Fable~5.2, Lean~4 formalisation of the results of~\cite{KNT},
\url{https://github.com/JWKNT/erdos1016}.

\end{thebibliography}

\end{document}
